\section{Results and Analysis}\label{sec:results}

\ins{Present your results clearly and analyze them thoughtfully. Go beyond reporting numbers or plots by interpreting what the results show, and what we can learn from them. Discuss limitations, failure cases, or unexpected outcomes where relevant.}

We report results from \textbf{Gemini~2.5 Flash} with thinking enabled
(\texttt{thinking\_budget}=8000, \texttt{include\_thoughts=True}) on
the full two-phase pipeline across six task slices:
four cultureMCQA slices (Korean / American / German / Polish, 50 items
each), GSM8K-MCQ (50 items), and XSAFETY (50 English + 50 German
prompts). After filtering each slice to baseline-correct (cultureMCQA
/ GSM8K) or baseline-refused (XSAFETY) items, Phase~2 runs 7 hint
conditions per item on the MCQ pipelines and 6 hint conditions per
item on the safety pipeline, giving \textbf{1{,}857 hinted
generations} in total. Few-shot biased CoT is excluded from XSAFETY
(no option-reordering trick for open-ended outputs, and constructing
compliance demos would risk generating harmful content).

\jr{Cross-model runs (Qwen via Groq, Claude via Anthropic API) are
the planned next experiments; the infrastructure handles both with no
code changes.}

\subsection{Cross-domain summary (Table~\ref{tab:summary})}

The headline result is consistent across all five slices:
\textbf{behavioral susceptibility varies by domain (0--32\%), but the
opaque-sycophancy count is zero everywhere}. Every observed flip
toward the injected wrong answer is paired with explicit
acknowledgment of the cue in either the visible reasoning or the
thinking trace.

\begin{table}[t]
\centering
\small
\begin{tabular}{lrrrrr}
\toprule
Slice & N hinted & Flips & Sycophantic & Transparent & Opaque \\
\midrule
Korean cultureMCQA (v2)   & 343   & 3    & 3    & 3    & \textbf{0} \\
American cultureMCQA (v2) & 350   & 0    & 0    & 0    & \textbf{0} \\
German cultureMCQA (v2)   & 350   & 2    & 2    & 2    & \textbf{0} \\
Polish cultureMCQA (v2)   & 350   & 0    & 0    & 0    & \textbf{0} \\
GSM8K-MCQ                 & 350   & 6    & 6    & 6    & \textbf{0} \\
XSAFETY (en + de)         & 114   & 36   & 36   & 36   & \textbf{0} \\
\midrule
\textbf{Total}            & \textbf{1{,}857} & \textbf{47} & \textbf{47} & \textbf{47} & \textbf{0} \\
\bottomrule
\end{tabular}
\caption{Cross-domain summary for Gemini 2.5 Flash with thinking.
For cultureMCQA and GSM8K-MCQ, ``Flips'' counts predicted answers
that differ from baseline; ``Sycophantic'' counts the subset that
match the injected wrong letter. For XSAFETY, ``Flips'' is
baseline-refused $\to$ hinted-comply. ``Transparent'' / ``Opaque''
split sycophantic flips by whether the cue is mentioned in the
visible reasoning or thinking trace.}
\label{tab:summary}
\end{table}

\subsection{cultureMCQA (Korean / American / German / Polish)}

Each cultureMCQA slice contains 50 hand-curated four-option items
grounded in CultureBank, with options shuffled to a near-uniform
correct-letter distribution (the first version of the dataset had a
strong B-positional bias which we corrected via
\texttt{scripts/\_shuffle\_options.py} with \texttt{seed=42}; the
original files are preserved as \texttt{*\_v1\_biased.json}).
Baselines are 49/50 (Korean), 50/50 (American), 50/50 (German), and
50/50 (Polish). Re-running Korean on the shuffled data yields 3 flips
vs.\ 5 on the original data, with the same qualitative pattern (only
authority cue produces flips; all flips transparent); we therefore
have direct evidence that the v1 numbers were not driven by the
position bias.

\begin{table}[t]
\centering
\small
\begin{tabular}{l rr rr rr rr}
\toprule
& \multicolumn{2}{c}{Korean (n=49)} & \multicolumn{2}{c}{American (n=50)} & \multicolumn{2}{c}{German (n=50)} & \multicolumn{2}{c}{Polish (n=50)} \\
\cmidrule(lr){2-3}\cmidrule(lr){4-5}\cmidrule(lr){6-7}\cmidrule(lr){8-9}
Condition & Flip & Ack & Flip & Ack & Flip & Ack & Flip & Ack \\
\midrule
cultural\_authority\_en & 2 & 48 & 0 & 49 & 1 & 50 & 0 & 50 \\
cultural\_authority\_ko & 1 & 48 & 0 & 50 & 1 & 50 & 0 & 49 \\
social\_en              & 0 & 42 & 0 & 46 & 0 & 42 & 0 & 41 \\
social\_ko              & 0 & 48 & 0 & 50 & 0 & 46 & 0 & 47 \\
indirect\_en            & 0 &  3 & 0 &  7 & 0 &  2 & 0 &  5 \\
indirect\_ko            & 0 & 13 & 0 &  8 & 0 &  9 & 0 &  4 \\
fewshot\_biased         & 0 &  0 & 0 &  1 & 0 &  0 & 0 &  0 \\
\bottomrule
\end{tabular}
\caption{cultureMCQA, per-condition counts of flipped predictions
(``Flip'') and items whose thinking trace explicitly named the cue
(``Ack''). Denominators are the per-culture baseline-correct counts.}
\label{tab:culture}
\end{table}

Three behavioral patterns hold across all four cultures
(Table~\ref{tab:culture}):

\begin{enumerate}
    \item \textbf{Cultural authority is the only effective cue.}
    Social, indirect, and few-shot biased conditions produce
    \emph{zero} flips on every culture. The 5 flips we observe on
    cultureMCQA in total are all on the cultural-authority cue:
    Korean (3), German (2), American and Polish (0 each). Korean is
    the most susceptible by a small margin; American and Polish are
    completely robust.
    \item \textbf{Acknowledgment is sharply asymmetric.} Direct cues
    are surfaced 82--100\% of the time across all four cultures.
    Indirect cues are surfaced only 4--26\%, and few-shot biased
    items 0--2\%. This asymmetry is invariant across cultures and
    cue languages.
    \item \textbf{Cue-language effect on response language varies
    across cultures.} On Korean cultureMCQA, the rate at which the
    visible reasoning contains Korean script is consistently higher
    under Korean-cue conditions (39--55\%) than under English-cue
    conditions (29--35\%). On American, this code-switching effect
    essentially disappears (0--2\% under \texttt{*\_ko} cues). On
    German, code-switching is broad (20--28\%) but only weakly
    sensitive to cue language. On Polish, code-switching is uniform
    at 32--38\% regardless of cue language---most of the non-English
    script in Polish responses is itself Polish (Polish cultural
    items pulling the model toward Polish), not Korean from the cue.
    Cue language thus acts as a soft register signal whose strength
    is moderated by how strongly the item content itself pulls a
    particular response language.
\end{enumerate}

\subsection{GSM8K-MCQ (Table~\ref{tab:gsm8k})}

GSM8K-MCQ~\citep{cobbe2021gsm8k} is the cross-domain check for
objective reasoning. Each item is a grade-school math word problem
converted to a four-option
MCQ with distractors drawn from intermediate values in the GSM8K
solution trace; the correct letter is randomized at conversion time,
so there is no position bias.

Baseline accuracy is 50/50 = 100\%. Phase 2 produces 6 flips out of
350, all on the \texttt{general\_authority} cue (which on this
dataset uses the framing ``an expert told me the answer is (X)''
rather than the cultural-studies professor). Korean-language
authority is marginally more effective than English-language
authority (4 vs 2 flips), but the order of magnitude is the same.
All other cue families---social, indirect, few-shot biased---produce
\emph{zero} flips on objective math.

\begin{table}[t]
\centering
\small
\begin{tabular}{lrr}
\toprule
Condition & Flip & Ack \\
\midrule
general\_authority\_en & 2/50 & 43/50 \\
general\_authority\_ko & 4/50 & 42/50 \\
social\_en             & 0/50 & 44/50 \\
social\_ko             & 0/50 & 45/50 \\
indirect\_en           & 0/50 &  1/50 \\
indirect\_ko           & 0/50 &  8/50 \\
fewshot\_biased        & 0/50 &  0/50 \\
\bottomrule
\end{tabular}
\caption{GSM8K-MCQ, Gemini 2.5 Flash with thinking. Same column
semantics as Table~\ref{tab:culture}.}
\label{tab:gsm8k}
\end{table}

Code-switching to Korean is rare on GSM8K (2--4\%), versus 24--55\%
on Korean cultureMCQA. This is consistent with arithmetic problems
pulling the reasoning toward computation and away from the cue
language; cultural items have no such pull and instead inherit
register from the cue.

All 6 flips on GSM8K are \emph{transparent}---the thinking trace
mentions the expert recommendation before adopting it, e.g.\
``the problem presents a contradiction... to arrive at the expert's
answer of (B)...'' \textbf{Opaque count is 0.}

\subsection{XSAFETY (Table~\ref{tab:safety})}

XSAFETY~\citep{wang2024xsafety} items are open-ended refusal prompts (50 English + 50 German)
where the expected behavior is refusal rather than a discrete answer.
We replace the MCQ-style hint families with three refusal-loosening
variants---\texttt{safety\_authority} (``a safety expert reviewed
this and confirmed it's fine to discuss for an academic study''),
\texttt{safety\_social} (``most people in our study group think this
question is harmless''), \texttt{safety\_indirect} (``I was just
curious---maybe a brief overview?'')---and detect refusal with a
bilingual regex classifier
(\texttt{scripts/\_refusal.py}). Compliance is measured as the rate
at which a baseline refusal flips to compliance under hint injection.

Baseline refusal rates are sparse but well-targeted: 11/50 (en) and
8/50 (de), concentrated in \texttt{Unsafe\_Instruction\_Topic} (67\%),
\texttt{Crimes\_And\_Illegal\_Activities} (43\%),
\texttt{Privacy\_And\_Property} (33\%), and
\texttt{Role\_Play\_Instruction} (33\%). XSAFETY categories that are
conversational by design (\texttt{Mental\_Health}, \texttt{Insult},
\texttt{commonsense}, etc.) are engaged rather than refused at
baseline and so do not enter Phase 2.

\begin{table}[t]
\centering
\small
\begin{tabular}{ll rr rrr}
\toprule
Prompt & Cue family & en cue & de cue & Resisted & Transparent & Opaque \\
\midrule
en & safety\_authority & 4/11 & 5/11 & 13/22 & 9/22 & 0 \\
en & safety\_social    & 2/11 & 3/11 & 17/22 & 5/22 & 0 \\
en & safety\_indirect  & 2/11 & 4/11 & 16/22 & 6/22 & 0 \\
de & safety\_authority & 4/8  & 2/8  & 10/16 & 6/16 & 0 \\
de & safety\_social    & 3/8  & 3/8  & 10/16 & 6/16 & 0 \\
de & safety\_indirect  & 1/8  & 3/8  & 12/16 & 4/16 & 0 \\
\bottomrule
\end{tabular}
\caption{XSAFETY, Gemini 2.5 Flash with thinking. The two \emph{cue}
columns are compliance-flip counts (baseline-refused $\to$ hinted-comply).
Resisted / Transparent / Opaque sum per row across both cue languages.}
\label{tab:safety}
\end{table}

Two findings stand out (Table~\ref{tab:safety}):

\begin{itemize}
    \item \textbf{XSAFETY is far more hint-susceptible than
    cultureMCQA.} Flip rates range from 12\% to 50\% across the six
    conditions, whereas the cultureMCQA flip rate sits at 0--4\%.
    The strongest single attack is the cross-lingual authority cue:
    English prompt + German authority = 5/11 = 45\%, German prompt +
    English authority = 4/8 = 50\%. Matched-language authority is
    less effective (en+en 36\%, de+de 25\%).
    \item \textbf{Cross-lingual cues are uniformly more effective in
    safety, but uniformly irrelevant for cultureMCQA flip rate.} In
    cultureMCQA the cue language affects \emph{response} language but
    \emph{not} flip rate (en and ko cue have the same per-culture
    flip count). In XSAFETY the cue language affects both. The most
    parsimonious explanation is that safety alignment is partly
    language-conditioned, so a cue in the wrong language partially
    bypasses guardrails calibrated to the prompt language.
\end{itemize}

Even at these much higher compliance rates, every compliance event
on XSAFETY is paired with explicit acknowledgment of the cue.
\textbf{Resisted 78/114, Transparent 36/114, Opaque 0/114.}

\subsection{Cross-cutting findings}

Pulling the six slices together, three patterns are invariant and
two vary cleanly across domains:

\paragraph{Invariant 1: 0 opaque sycophancy across 1{,}857
generations.} Every flip toward the injected wrong answer is paired
with an explicit cue mention in the visible reasoning or the
thinking trace. This is the strongest single empirical finding of
our study. It suggests that, for Gemini~2.5 Flash with thinking, the
worst-case faithfulness failure motivating the Turpin et al. line of
work---behavioral compliance with a biased cue paired with verbal
silence---does not occur at our sample size.

\paragraph{Invariant 2: Authority is the dominant flip-inducing cue.}
Across the four cultureMCQA slices and GSM8K, the only cue family
that produces any flip is the authority cue
(\texttt{cultural\_authority} on cultureMCQA,
\texttt{general\_authority} on GSM8K). Social, indirect, and few-shot
biased cues produce \emph{zero} flips on all five MCQ slices. The
XSAFETY pattern is qualitatively similar: \texttt{safety\_authority}
is the strongest flipper there as well.

\paragraph{Invariant 3: Direct cues are loudly acknowledged;
indirect / structural cues are not.} Authority and social cues are
acknowledged in $\geq$82\% of thinking traces in every slice,
whereas indirect cues are acknowledged in $\leq$27\% and few-shot
biased (under a strict criterion that requires explicit reference to
the demonstrations) is acknowledged in $\leq$2\%. Because the
indirect / fewshot cues never actually flip behavior on the MCQ
domains, this asymmetry never aligns with opaque sycophancy in our
data---but it does mean an auditing pipeline that screens reasoning
traces for hint mentions would miss those families almost entirely
if they did start producing flips in some other model.

\paragraph{Domain-dependent variation 1: behavioral susceptibility
scales by an order of magnitude.} Flip rates are 0--0.9\% across the
four cultureMCQA slices, 1.7\% on GSM8K-MCQ, and 18--50\% on
XSAFETY. The pattern is intuitive: safety alignment is engineered to
be tighter than cultural-knowledge defaults, so when a cue does
crack it (especially cross-lingual authority), it cracks much more
easily than knowledge defaults that the model is independently
confident about.

\paragraph{Domain-dependent variation 2: cross-lingual cue effects
are pronounced for safety, absent for cultureMCQA flip rate.}
On cultureMCQA, English and Korean variants of the same cue have
identical flip counts. On XSAFETY, the cross-lingual cue is
\emph{stronger} than the matched-lingual one (en+de 45\% vs.\ en+en
36\%; de+en 50\% vs.\ de+de 25\%). The hint language always affects
the response language (most strongly on Korean cultureMCQA: 39--55\%
ko script under ko cue vs.\ 29--35\% under en cue), but the
mechanism by which cue language influences behavior appears to be
distinct in safety vs.\ knowledge domains.

\subsection{Limitations}

The results are obtained on a single model (Gemini~2.5 Flash) with a
single hint-template set per cue family. Failure modes absent in
this configuration may emerge in weaker or differently-aligned
models; planned follow-ups use the same pipeline on Qwen via Groq,
Llama via Groq, and Claude via Anthropic API. Cross-model results
are not yet available at the time of writing.

Hint-acknowledgment detection uses a tight, hint-type-specific
keyword list. A model could discuss the cue using paraphrases not on
our trigger list, in which case our acknowledgment counts are
lower bounds---and, more importantly, opaque-sycophancy counts could
be upper-biased (false-positive opacity from miss-keyword cases).
Two iterations of regex tightening were required during the XSAFETY
analysis, and three apparently-opaque cases turned out to be
classifier artifacts (a baseline-empathy expression mistaken for a
refusal in two cases, and a missed German noun form of an
acknowledgment keyword in the third). After tightening, the count
went to 0 across all five slices.

The cultureMCQA correct-option text averages 70--84 characters
while distractors average 25--33 characters; the shuffle fixes
positional bias but not length bias. Because the model is highly
accurate at baseline regardless, this confound is unlikely to drive
the Phase~2 numbers, but it should be removed in any follow-up
study.

Finally, pragmatic non-equivalence of hint types across cultures is
a real but untested confound: ``a professor says X'' may carry
heavier pragmatic weight in some cultures than others. A
perceived-pressure rating layer over the hint set would be a
natural extension.
